\documentclass[11pt,a4paper]{article}
\usepackage[utf8]{inputenc}
\usepackage[T1]{fontenc}
\usepackage[spanish,es-tabla]{babel}
\usepackage{lmodern,microtype}
\usepackage[a4paper,margin=2.25cm,headheight=14pt]{geometry}
\usepackage{amsmath,amssymb,booktabs,longtable,array,graphicx,xcolor}
\usepackage{caption,float,placeins,fancyhdr,enumitem}
\usepackage{hyperref}
\definecolor{azul}{HTML}{20566F}
\hypersetup{hypertexnames=false,colorlinks=true,linkcolor=azul,urlcolor=azul,citecolor=azul,pdfauthor={Pablo Sánchez González},pdftitle={La composición del gasto público antes y después de la pandemia}}
\setlength{\parindent}{1em}
\setlength{\parskip}{3pt}
\setlength{\emergencystretch}{2em}
\linespread{1.05}
\captionsetup{font=small,labelfont=bf,justification=raggedright,singlelinecheck=false}
\pagestyle{fancy}\fancyhf{}
\fancyhead[L]{\footnotesize Composición institucional del gasto público}
\fancyhead[R]{\footnotesize UE-27, 2015--2024}
\fancyfoot[C]{\thepage}
\renewcommand{\headrulewidth}{0.3pt}
\newcommand{\nota}[1]{\par\vspace{2pt}{\footnotesize\noindent #1\par}}
\newcommand{\autor}{Pablo Sánchez González}
\newcommand{\afiliacion}{Universidad Complutense de Madrid \textperiodcentered{} Economía Pública}
\makeatletter\let\tableinput\@@input\makeatother
\addto\captionsspanish{\renewcommand{\refname}{Fuentes de datos y documentación}}
\input{cifras.tex}

\begin{document}
\begin{titlepage}
\thispagestyle{empty}
\vspace*{1.5cm}
\begin{center}
{\LARGE\bfseries La composición del gasto público\par}
\vspace{6pt}
{\LARGE\bfseries antes y después de la pandemia\par}
\vspace{10pt}
{\large Evidencia de los países de la Unión Europea, 2015--2024\par}
\vspace{1cm}
{\large\autor\par}
\vspace{6pt}
Universidad Complutense de Madrid\\
Economía Pública
\vspace{12pt}\par
24 de septiembre de 2026
\end{center}
\vspace{1cm}
\begin{abstract}
La pandemia sometió a las administraciones públicas a una fuerte presión para sostener las rentas y financiar los servicios públicos. Este trabajo estudia si, después de los primeros años de la crisis, cambió el peso relativo de los distintos niveles de gobierno en el gasto público de los países de la Unión Europea. A partir de datos de Eurostat, se compara la distribución del gasto entre los gobiernos centrales, regionales y locales y los fondos de seguridad social en 2017--2019 y 2022--2024. En el promedio de los 27 países, la participación de los gobiernos centrales aumenta aproximadamente 1,15 puntos porcentuales, mientras disminuye la de las administraciones locales y la de los fondos de seguridad social. La mayor participación del gobierno central se observa en 21 países, aunque la intensidad y la dirección de los cambios difieren entre economías. El resultado se mantiene al modificar los periodos de comparación y al excluir países individualmente. La evidencia apunta a un cambio moderado en la composición institucional del gasto, compatible con un mayor papel de los gobiernos centrales en su financiación. Esta evolución no permite atribuir los cambios exclusivamente a la pandemia ni concluir que se haya producido una transferencia equivalente de competencias desde las administraciones territoriales.
\end{abstract}
\vspace{12pt}
\noindent\textbf{Palabras clave:} gasto público; relaciones fiscales intergubernamentales; composición del gasto; COVID-19; Unión Europea.\\[4pt]
\textbf{Clasificación JEL:} H50, H70, H77.
\vfill
\end{titlepage}
\setcounter{page}{1}

\section{Introducción}
La respuesta a la pandemia exigió un esfuerzo extraordinario de las administraciones públicas. Además de atender nuevas necesidades sanitarias, los gobiernos tuvieron que sostener las rentas de hogares y empresas y asegurar la continuidad de los servicios públicos. Estas actuaciones plantean una cuestión que va más allá del volumen del gasto: cómo se distribuyó ese esfuerzo entre los distintos niveles de gobierno y si la distribución observada después de los primeros años de la crisis difiere de la anterior.

Este trabajo estudia los cambios en la composición institucional del gasto público en los 27 países de la Unión Europea. Se compara la participación de los gobiernos centrales, regionales y locales y de los fondos de seguridad social durante 2017--2019 y 2022--2024. El análisis de la trayectoria anual entre 2015 y 2024 permite situar esa comparación en un contexto temporal más amplio y distinguir los primeros años de la pandemia de su evolución posterior.

La pregunta es relevante para la economía pública porque las responsabilidades de gasto y la capacidad para financiarlas se distribuyen de manera desigual entre administraciones. Ante una perturbación que afecta simultáneamente a todo el territorio, los gobiernos centrales pueden asumir una mayor parte de la financiación, incluso cuando los servicios siguen prestándose desde las administraciones regionales o locales. Observar un cambio en las participaciones de gasto ayuda a describir esa evolución, aunque no basta para conocer cómo se modificaron las competencias o la autonomía presupuestaria.

Los resultados muestran una mayor participación de los gobiernos centrales en el gasto: el incremento medio es de aproximadamente 1,15 puntos porcentuales y se observa en 21 países. A su vez, disminuye el peso medio de las administraciones locales y de los fondos de seguridad social. Este patrón presenta diferencias nacionales importantes y su magnitud depende del periodo de referencia. La contribución del trabajo consiste en documentar esos cambios, valorar su extensión y examinar hasta qué punto se mantienen bajo otras comparaciones. Su alcance es descriptivo: la coincidencia temporal con la pandemia no permite aislar su efecto de otras decisiones y perturbaciones económicas.

\section{La distribución del gasto entre niveles de gobierno}
\subsection{Estabilización y prestación de servicios públicos}
La distribución del gasto entre administraciones depende de las funciones que desempeña cada una y de sus relaciones financieras. La prestación de servicios territoriales, el apoyo a las rentas y la financiación de actuaciones comunes pueden recaer en niveles de gobierno diferentes. Para interpretar los cambios en sus participaciones, es necesario distinguir quién financia una actuación, quién la ejecuta y dónde se registra el gasto.

Desde esta perspectiva, cabe esperar que una crisis de alcance nacional modifique el reparto del esfuerzo presupuestario. El gobierno central puede financiar programas de apoyo a las rentas o aportar recursos a otras administraciones. Por su parte, el gasto de los gobiernos territoriales dependerá de los servicios que tengan encomendados y de los recursos de que dispongan. Cuando las prestaciones se canalizan a través de fondos de seguridad social, estos también pueden adquirir un mayor protagonismo. No existe, por tanto, una respuesta uniforme que pueda deducirse al margen de la organización institucional de cada país.

Una mayor participación del gobierno central puede ser compatible tanto con un aumento de su gasto directo como con mayores transferencias a otras administraciones. Estas son posibles vías de interpretación, no mecanismos demostrados por los resultados. El análisis examina la distribución contable del gasto y no permite identificar por separado la contribución de cada vía.

\subsection{Participación en el gasto y autonomía fiscal}
El peso presupuestario de una administración y su capacidad de decisión son dimensiones relacionadas, pero distintas. Una administración local puede aumentar su gasto y, aun así, perder participación en el conjunto si el gasto de otras administraciones crece más deprisa. Del mismo modo, una mayor financiación procedente del gobierno central no implica necesariamente que este asuma la gestión de los servicios.

Las transferencias entre administraciones son especialmente importantes para interpretar las cifras. Si el gobierno central transfiere diez unidades a una administración local y esta las emplea en servicios, la suma del gasto de ambas puede recoger tanto la transferencia como su utilización. El gasto consolidado elimina las operaciones internas correspondientes. Por ello, las participaciones calculadas sobre la suma del gasto de los distintos niveles describen su distribución institucional, pero no un reparto exclusivo del gasto final.

Estas distinciones delimitan la interpretación económica. Una mayor participación de los gobiernos centrales no demuestra por sí misma una mejora de la eficiencia ni una pérdida de autonomía territorial. Para valorar esas consecuencias sería necesario estudiar también las competencias, las condiciones de financiación, los costes y los resultados de los servicios públicos.

\section{Datos y medición}
\subsection{Fuente y ámbito del estudio}
Se utiliza la base de Eurostat \texttt{gov\_10a\_exp}, gasto total y todas las funciones COFOG, con frecuencia anual \cite{datos}. Se emplea la versión de los datos actualizada el 16 de septiembre de 2026, que permite seguir a los mismos 27 países durante diez años, de 2015 a 2024. La base reúne, por tanto, 270 observaciones anuales de países.

Los cuatro subsectores son gobierno central (S.1311), estatal o regional (S.1312), local (S.1313) y fondos de seguridad social (S.1314). S.13 se conserva como referencia consolidada. ``Regional'' designa aquí el subsector S.1312, no cualquier institución territorial que lleve ese nombre; ``seguridad social'' tampoco equivale a todo el gasto funcional en protección social.

\subsection{Diferencias institucionales y comparabilidad}
La clasificación estadística no presenta los mismos subsectores por separado en todos los países. En la Unión Europea, el subsector de gobierno estatal o regional S.1312 se identifica en Austria, Bélgica, Alemania y España \cite{sectores}. En Malta, las operaciones sociales pertinentes se integran en las cuentas del gobierno central y no se presenta un subsector separado de fondos de seguridad social \cite{malta}.

Estas diferencias explican las 230 observaciones ausentes del subsector regional y las diez de seguridad social de Malta. Para construir una composición comparable de cuatro componentes se les asigna un valor cero, limitado a esos casos documentados. La decisión se refiere a la clasificación de las administraciones: no supone que los demás países carezcan de instituciones regionales o que Malta no tenga prestaciones sociales. La justificación detallada y la conservación de las marcas originales se recogen en el cuaderno metodológico.

También deben considerarse las revisiones de las cuentas y los cambios de registro. Los metadatos de Eurostat señalan cuestiones de comparabilidad en Finlandia, Eslovaquia e Irlanda \cite{meta}. En este último caso, la publicación del subsector de seguridad social incorpora una serie retrospectiva; no debe interpretarse automáticamente como un cambio económico ocurrido en el año de publicación. Las comprobaciones posteriores valoran la influencia de estos países en los resultados.

\subsection{Medición de la composición del gasto}
Para el país $i$, subsector $k$ y año $t$, sea $E_{ikt}$ el gasto publicado. La participación de cada administración en la suma del gasto de los cuatro subsectores se define como
\begin{equation}
 D_{it}=\sum_{k=1}^{4} E_{ikt},\qquad
 s_{ikt}=100\frac{E_{ikt}}{D_{it}},\qquad \sum_k s_{ikt}=100.
\label{eq:share}
\end{equation}
La consolidación explica que $D_{it}$ no coincida con S.13 \cite{meta}. En el panel, la diferencia representa en promedio 11,025 pp del PIB, con valores entre 0,2 y 25,3. La diferencia responde al distinto tratamiento de las operaciones entre administraciones, por lo que no puede interpretarse como gasto adicional en servicios públicos.

La especificación principal emplea cifras expresadas en porcentaje del PIB, para mantener una unidad común de presentación. El PIB común se cancela en la razón de la ecuación~\eqref{eq:share}; la única diferencia frente a calcular con niveles nominales procede del redondeo y de la precisión publicada. Se verifica este punto con los valores en millones de euros publicados en la misma versión de los datos.

Se asigna igual peso a los países. De este modo, el resultado resume la experiencia de los países de la UE-27 sin que las economías de mayor tamaño dominen el promedio. No representa la participación de cada nivel de gobierno en el gasto agregado europeo, que requeriría una ponderación basada en importes monetarios.

\section{Estrategia empírica}
\subsection{Comparación entre la etapa anterior y la posterior a la crisis}
La comparación principal utiliza dos medias de tres años:
\begin{equation}
 \Delta_{ik}=\frac{1}{3}\sum_{t=2022}^{2024}s_{ikt}
             -\frac{1}{3}\sum_{t=2017}^{2019}s_{ikt},
 \qquad \overline\Delta_k=\frac{1}{27}\sum_i\Delta_{ik}.
\end{equation}
La pregunta es si la composición media de una etapa posterior difiere de la etapa inmediatamente anterior a la crisis. Promediar amortigua la influencia de un año singular y evita que 2020 determine por sí solo la comparación. Se utiliza como referencia el trienio inmediatamente anterior a la pandemia y se reserva 2020--2021 para examinar sus primeros años. La elección de periodos forma parte del diseño del estudio y se contrasta con otras alternativas.

Los periodos no reemplazan el examen temporal. La agregación pierde información sobre cuándo ocurre el cambio y puede ocultar reversión, tendencias o reformas. Por eso se muestra la serie anual completa, se distingue 2020--2021 de 2022--2024 y se examinan otras ventanas. La media de cuotas tampoco equivale a la cuota del gasto acumulado: la primera pondera igual los años; la segunda los ponderaría según el gasto total.

Una regresión de serie temporal sería útil para otra pregunta, pero diez observaciones anuales, cinco anteriores a 2020, ofrecen información limitada para separar tendencia, salto y persistencia. Un panel no elimina el problema de identificación: un indicador post común es colineal con efectos fijos completos de año; sin ellos puede absorber shocks simultáneos. La comparación entre periodos permite responder de forma directa a la pregunta sobre la distribución media del gasto, mientras que la evolución anual ayuda a interpretar cuándo se produjeron los cambios. Ninguno de los dos enfoques, por sí solo, permite atribuirlos causalmente a la pandemia.

\subsection{Diferencias nacionales y evaluación estadística}
Para describir la intensidad de la recomposición se calcula
\begin{equation}
 R_i=\frac12\sum_k|\Delta_{ik}|.
\end{equation}
Como las diferencias suman cero, el índice equivale a la suma de aumentos y toma valores de 0 a 100 pp. No identifica una transferencia presupuestaria observada. Se utiliza como medida descriptiva de la intensidad del cambio, sin atribuirle una interpretación causal.

Los tests utilizan 27 diferencias nacionales, no 270 observaciones independientes. La referencia principal es el test $t$ bilateral de media cero, con intervalos marginales del 95\,\% y ajuste de Holm para una familia de tres contrastes: central, local y seguridad social. El nivel regional se conserva en la composición, pero se describe sin una prueba sustentada en 23 ceros estructurales. Wilcoxon se muestra como análisis complementario, con otra familia de tres ajustes, para comprobar si las conclusiones dependen del procedimiento estadístico.

Los países no son una muestra aleatoria de una población mayor: constituyen todo el ámbito elegido. La inferencia requiere interpretar las diferencias como realizaciones de un modelo más general y asumir independencia entre países; la referencia exacta del test $t$ requiere además normalidad. Los shocks comunes y la dependencia composicional limitan esa interpretación. Holm corrige multiplicidad, no esos problemas. En consecuencia, la interpretación combina la magnitud de los cambios, su distribución entre países y las comprobaciones de robustez.

\section{Resultados}
\subsection{Los gobiernos centrales ganan peso en la distribución del gasto}
La participación media de los gobiernos centrales pasa del \datoCentralPre{}\,\% al \datoCentralPost{}\,\% entre los dos trienios, lo que supone un incremento de \datoCentral{} puntos porcentuales (pp). En paralelo, las administraciones locales pierden alrededor de 0,46 puntos de participación y los fondos de seguridad social, 0,67. El peso del subsector regional apenas varía en el promedio europeo, que incluye los países donde ese nivel no se presenta por separado.
\begin{table}[H]\centering\small
\caption{Composición media por etapa: media simple de 27 países}
\begin{tabular}{@{}lrrrr@{}}\toprule
Subsector & 2017--2019 & 2020--2021 & 2022--2024 & Cambio (pp)\\\midrule
\tableinput{tablas/periodos.tex}
\bottomrule\end{tabular}
\nota{Fuente: elaboración propia con Eurostat. Participaciones sobre la suma del gasto de los cuatro subsectores. Cambio: 2022--2024 menos 2017--2019.}
\end{table}
La evolución anual de la figura~\ref{fig:serie} ayuda a situar estos cambios. El peso de los gobiernos centrales había disminuido entre 2015 y 2019, pero aumenta en 2020 y, durante 2022--2024, se mantiene en promedio cerca del nivel de los dos primeros años de la pandemia. Las administraciones locales recuperan entonces parte de la participación perdida, mientras que los fondos de seguridad social reducen su peso. La composición del gasto continúa cambiando después de 2021, aunque la participación media de los gobiernos centrales apenas se modifique entre ambas etapas.
\begin{figure}[H]\centering
\includegraphics[width=.95\textwidth]{figuras/serie_anual.pdf}
\caption{Participación media de cada nivel de gobierno en el gasto, UE-27, 2015--2024. La franja identifica 2020--2021; cada panel tiene su propia escala vertical.}\label{fig:serie}
\nota{Fuente: elaboración propia con Eurostat. Media simple de 27 países; porcentaje de la suma del gasto de cuatro subsectores.}
\end{figure}

\subsection{Alcance de la evidencia estadística}
La evidencia estadística es más consistente para el incremento de la participación de los gobiernos centrales. Su intervalo de confianza del 95\,\% va de \datoCentralLow{} a \datoCentralHigh{} pp, y el contraste de medias mantiene un valor $p$ de \datoPHolm{} después del ajuste de Holm. Las reducciones de la participación local y de seguridad social no resultan significativas al 5\,\% con ese mismo criterio. El contraste de Wilcoxon sí encuentra diferencias en los tres casos, aunque evalúa propiedades distintas de la distribución. Por ello, la evidencia sobre esas reducciones debe interpretarse con mayor cautela.
\begin{table}[H]\centering\small
\caption{Cambios nacionales e inferencia exploratoria}
\begin{tabular}{@{}lrrrr@{}}\toprule
 & Media (pp) & IC 95\,\% & $p_{t,\mathrm{Holm}}$ & $p_{W,\mathrm{Holm}}$\\\midrule
\tableinput{tablas/contrastes.tex}
\bottomrule\end{tabular}
\nota{27 países. Intervalos marginales sin ajuste; Holm por separado para tres tests $t$ y tres Wilcoxon. No se controla una selección posterior entre las seis pruebas.}
\end{table}
\begin{figure}[H]\centering
\includegraphics[width=.92\textwidth]{figuras/intervalos.pdf}
\caption{Diferencias medias 2022--2024 frente a 2017--2019 e intervalos marginales.}
\nota{Fuente: elaboración propia. La referencia vertical es cero; los IC no son simultáneos ni están ajustados por Holm.}
\end{figure}
El intervalo de la variación de la participación local puede excluir cero y su valor $p$ ajustado ser mayor que 0,05 porque responden a criterios distintos. La discrepancia de seguridad social entre procedimientos debe leerse junto a la heterogeneidad: muchas caídas coexisten con una subida grande en Chipre. Ningún intervalo incorpora incertidumbre por revisiones, clasificación o elección de ventanas.

\subsection{Un patrón compartido con diferencias importantes entre países}
La participación del gobierno central en el gasto aumenta en 21 países y disminuye en seis. Los mayores incrementos se observan en Países Bajos, Hungría y Letonia, donde superan los 3,5 pp. En España, el aumento es de \datoEspana{} pp y se acompaña de una menor participación de los gobiernos regionales y locales y de los fondos de seguridad social. Estas reducciones son relativas: no implican necesariamente que su gasto haya disminuido en términos monetarios.

\begin{figure}[H]\centering
\includegraphics[width=.97\textwidth,height=.70\textheight,keepaspectratio]{figuras/comparacion_paises.pdf}
\caption{Cambio en la participación del gobierno central en el gasto por país.}\label{fig:central-paises}
\nota{Fuente: elaboración propia con Eurostat. Diferencias entre 2022--2024 y 2017--2019, en puntos porcentuales. La línea vertical indica ausencia de cambio.}
\end{figure}

La figura~\ref{fig:paises} muestra que el reparto de los cambios entre administraciones varía considerablemente. En algunos países, la mayor participación del gobierno central coincide sobre todo con una pérdida de peso de la seguridad social; en otros, con una reducción del peso de los gobiernos locales. Chipre destaca por la evolución contraria: el gobierno central pierde 6,32 pp de participación y los fondos de seguridad social ganan 6,62. Este caso ilustra por qué la media europea no basta para describir la experiencia de cada país.
\begin{figure}[H]\centering
\includegraphics[width=.97\textwidth,height=.70\textheight,keepaspectratio]{figuras/mapa_cambios.pdf}
\caption{Cambios en la participación de cada nivel de gobierno en el gasto.}\label{fig:paises}
\nota{Fuente: elaboración propia con Eurostat. Diferencias entre 2022--2024 y 2017--2019, en puntos porcentuales. La escala es común a todos los países y subsectores.}
\end{figure}
El índice de reasignación resume la intensidad de estos movimientos sin atender a su dirección. Su media es de \datoReasignacion{} pp y su mediana, de 1,746. España presenta un valor de \datoEspana{} pp, inferior al promedio europeo. Un cambio de esta magnitud representa una modificación moderada de la distribución del gasto, aunque las diferencias entre países impiden describir un proceso uniforme.

\section{Robustez y límites de comparabilidad}
\subsection{La elección del periodo de comparación}
En las seis comparaciones entre etapas anteriores y posteriores a la pandemia de la tabla~\ref{tab:ventanas}, la participación media de los gobiernos centrales aumenta entre \datoVentanaMin{} y \datoVentanaMax{} pp. Si se amplía el periodo de referencia a 2015--2019, el incremento es de \datoBaseAmplia{} pp. En cambio, entre los bienios 2015--2016 y 2018--2019 esa participación había disminuido aproximadamente 0,82 pp. La dirección del cambio posterior se mantiene, pero su magnitud debe leerse teniendo en cuenta la evolución previa.
\begin{table}[H]\centering\small
\caption{Variación del peso de los gobiernos centrales según el periodo comparado}\label{tab:ventanas}
\begin{tabular}{@{}llllr@{}}\toprule
Comparación & Base & Etapa comparada & Media (pp) & Subidas\\\midrule
\tableinput{tablas/ventanas.tex}
\bottomrule\end{tabular}
\nota{Fuente: elaboración propia. 27 países en cada fila. Las seis primeras filas son comparaciones pre/post. Las restantes describen fases y una referencia precrisis, no placebos causales.}
\end{table}
Las comparaciones permiten valorar la importancia de la elección temporal. El uso de años individuales concede mayor influencia a acontecimientos puntuales, mientras que una referencia más larga incorpora una etapa anterior distinta. No hay dos trienios anteriores no solapados disponibles en 2015--2019; por eso la referencia precrisis utiliza bienios y su índice no se compara como si tuviera idéntico diseño.

\subsection{Países influyentes, revisiones y precisión}
Al repetir la comparación excluyendo un país cada vez, el incremento medio de la participación de los gobiernos centrales se mantiene entre \datoMinLoo{} y \datoMaxLoo{} pp. El resultado, por tanto, no depende de un único país. Al calcular las mismas participaciones con datos en millones de euros, se obtiene una variación de \datoEuroCentral{} pp, prácticamente idéntica a la principal.

La tabla~\ref{tab:calidad} examina Malta y tres casos con cautelas documentales. Al excluir conjuntamente Finlandia, Eslovaquia e Irlanda, el peso medio de los gobiernos centrales aumenta \datoCalidad{} pp. Estas comprobaciones permiten valorar la influencia de países con particularidades estadísticas, aunque modifican el conjunto estudiado y no eliminan todas las dificultades de comparación.
\begin{table}[H]\centering\small
\caption{Sensibilidad a la composición de la muestra}\label{tab:calidad}
\begin{tabular}{@{}lrrrr@{}}\toprule
Muestra & Países & Central & Local & Seg. social\\\midrule
\tableinput{tablas/calidad.tex}
\bottomrule\end{tabular}
\nota{Fuente: elaboración propia. Cambios medios en pp, 2022--2024 menos 2017--2019. Se repite la comparación con distintos conjuntos de países.}
\end{table}

\paragraph{Comprobación adicional de clasificación.}
Como ejercicio complementario, la suma de las participaciones del gobierno central y de la seguridad social aumenta \datoAgrupado{} pp. Esta agrupación ayuda a valorar diferencias de registro entre países, pero responde a una clasificación distinta de la empleada en el análisis principal. Los resultados desagregados siguen siendo la referencia para estudiar cómo cambia la composición del gasto.

\section{Discusión y conclusiones}
El gasto público de los países de la Unión Europea presenta una distribución institucional distinta en 2022--2024 de la observada en 2017--2019. Los gobiernos centrales ganan, en promedio, alrededor de 1,15 puntos porcentuales de participación, mientras que disminuye el peso de las administraciones locales y de los fondos de seguridad social. El cambio se observa en una amplia mayoría de países, aunque no sigue la misma intensidad ni se compensa de igual forma entre los restantes niveles de gobierno.

La evolución es compatible con un mayor protagonismo de los gobiernos centrales en la financiación del gasto después de los primeros años de la pandemia. Su capacidad para sostener rentas y apoyar a otras administraciones ofrece una interpretación económica plausible. Sin embargo, los datos utilizados no permiten determinar qué parte del cambio responde a transferencias, a gasto directamente ejecutado o a otras decisiones presupuestarias. Tampoco permiten atribuirlo exclusivamente a la crisis sanitaria, puesto que los años posteriores recogen otras perturbaciones y políticas.

La distinción entre participación en el gasto y autonomía fiscal resulta esencial. Una administración territorial puede perder peso relativo y continuar prestando los mismos servicios, e incluso aumentar su gasto. Por ello, los resultados no prueban un traslado equivalente de competencias hacia el gobierno central ni permiten evaluar si la nueva distribución es más eficiente.

La principal conclusión es que la composición del gasto cambió de forma moderada y con importantes diferencias nacionales. El mayor peso de los gobiernos centrales constituye el rasgo común más claro, pero el análisis de los demás niveles de gobierno es necesario para entender cómo se produjo esa transformación. El estudio de las transferencias intergubernamentales y de las funciones de gasto permitiría explicar mejor las distintas trayectorias observadas.

\paragraph{Datos y materiales complementarios.}
La entrega incluye los datos utilizados, los cuadernos de análisis y metodología y las instrucciones para reproducir tablas, figuras y documentos. El cuaderno metodológico desarrolla el tratamiento de los subsectores no presentados por separado y la elección de los periodos de comparación.

\begin{thebibliography}{9}
\small\raggedright
\bibitem{datos} Eurostat (2026). \emph{General government expenditure by function (COFOG)}, gov\_10a\_exp. Instantánea actualizada el 16 de septiembre de 2026. \url{https://doi.org/10.2908/GOV_10A_EXP}.
\bibitem{meta} Eurostat (2026). \emph{COFOG: reference metadata}, apartados 18.5 y 19. \url{https://webgate.ec.europa.eu/eurostat/cache/metadata/en/gov_10a_exp_esms.htm}.
\bibitem{sectores} Eurostat (2023). \emph{Glossary: General government sector}, p. 1. \url{https://ec.europa.eu/eurostat/statistics-explained/SEPDF/cache/1123.pdf}.
\bibitem{malta} National Statistics Office, Malta (2022). \emph{EDP Consolidated Inventory of Sources and Methods, ESA 2010}, apartados 1.2--1.4. \url{https://nso.gov.mt/wp-content/uploads/2023/01/EDP-Inventory-ESA-2010.pdf}.
\end{thebibliography}

\clearpage
\appendix
\section{Resultados nacionales}
\small
\begin{longtable}{@{}lrrrrrr@{}}
\caption{Cambios en pp y distancia composicional}\\
\toprule País & Central & Regional & Local & Seg. social & $R_i$ & Central+SS\\\midrule
\endfirsthead
\toprule País & Central & Regional & Local & Seg. social & $R_i$ & Central+SS\\\midrule
\endhead
\bottomrule\endfoot
\tableinput{tablas/paises.tex}
\end{longtable}
\nota{Fuente: elaboración propia. $R_i$ es no negativo; las variaciones de las cuatro cuotas suman cero antes del redondeo. Central+SS suma dos componentes con el denominador original, sin consolidación adicional.}
\normalsize
\end{document}
